\clearpage\thispagestyle{empty}
\pdfbookmark[0]{Quellen zu Komet (19. Dezember 1811)}{comet1811dec19-sources}
\begin{center}\textbf{Quellen zu Komet (19. Dezember 1811)}\end{center}
\noindent Diese KI-Übertragung ins heutige Deutsch folgt dem Bericht in \href{https://gdz.sub.uni-goettingen.de/id/PPN236007467}{\textit{Werke}, Band~VI} (Göttingen, 1874), S.~342--343. Der Erstdruck erschien unter der Überschrift Göttingen in den \textit{Göttingischen gelehrten Anzeigen}, Heft 201 vom 19. Dezember 1811, S.~\href{https://api.digitale-sammlungen.de/iiif/image/v2/bsb10538680_00977/full/1600,/0/default.jpg}{2001}, \href{https://api.digitale-sammlungen.de/iiif/image/v2/bsb10538680_00978/full/1600,/0/default.jpg}{2002}, \href{https://api.digitale-sammlungen.de/iiif/image/v2/bsb10538680_00979/full/1600,/0/default.jpg}{2003} und \href{https://api.digitale-sammlungen.de/iiif/image/v2/bsb10538680_00980/full/1600,/0/default.jpg}{2004}. Am Ende des Berichts steht das Datum 12. Dezember. Das digitalisierte Exemplar stammt aus der Bayerischen Staatsbibliothek.
\par\smallskip\noindent Beim Zeitpunkt des Periheldurchgangs stehen in der Zeitschrift auf S.~2002 35 Sekunden und in der Werkausgabe 55. Auf S.~2003 druckt die Zeitschrift Nicolais sechzehnzeilige Ephemeride für Januar 1812 ab. An dieser Stelle stehen in der Werkausgabe drei Striche. Auch die Form Blätten und die beiden Punkte in den Rektaszensionen für den 17. und 20. November in der Marseiller Tabelle folgen der Werkausgabe. Den beiden Punkten wird hier keine numerische Bedeutung zugewiesen. KI-Übertragung und KI-Satz.
\clearpage
